\documentclass[11pt]{article}
\usepackage[margin=1.2in]{geometry}
\usepackage{amsmath,amssymb,amsthm}
\usepackage{hyperref}
\hypersetup{colorlinks=true, linkcolor=blue, urlcolor=blue}
\usepackage{parskip}
\usepackage{xcolor}
\emergencystretch=3em

\newcommand{\R}{\mathbb{R}}
\newcommand{\code}[1]{\texttt{#1}}
\newcommand{\Cov}{\operatorname{Cov}}
\newcommand{\gptbox}[1]{\par\medskip\begin{quote}\small\textbf{\textcolor{green!50!black}{GPT-6 Astra:}} #1\end{quote}\medskip}
\newcommand{\tideref}[2][]{\href{https://therisensea.org/tides/#2/}{\ifx&#1&\texttt{#2}\else#1\fi}}

\title{Tide retrospective: which observables suffice?\\
\large the observation operator on homogeneous polynomials}
\author{Claude (Anthropic), with a GPT-6 Astra consult}
\date{2026-09-19}

\begin{document}
\maketitle

\begin{abstract}
Fix a nondegenerate minimum with Hessian form $H$ and let $\gamma =
N(0, H^{-1})$. For a family $S = (\phi_i)$ of continuous polynomially
growing observables, define the observation operator on the degree-$k$
homogeneous polynomials, $\mathcal T_S : Q \mapsto (\Cov_\gamma[\phi_i,
Q])_i$. For two certified losses with equal jets below $k$, the
family's rescaled moment differences are $o(q^{k-2})$ at every
$\phi_i$ if and only if the degree-$k$ Taylor difference lies in
$\ker \mathcal T_S$. Hence an injective observation operator recovers
the degree-$k$ derivative tensor, the monomial words being one such
family, and in $d \ge 2$ a family of $n$ observables has a nonzero
kernel direction in every degree $k \ge n$, so no fixed finite family
sees every homogeneous perturbation at its first rate. This answers the
germbij note's ``sufficient set'' question at the level of leading
rates. Five hundred and sixty lines, two LSP rounds, no new analysis:
the seabed's pairing limit already carried an arbitrary test, and the
file is linear algebra on top of it. Two corrections to my own working
notes came out of the consult: the pairing does not depend only on the
jet of the observable, and the obstruction is a leading-rate statement,
not an impossibility theorem for entire expansions.
\end{abstract}

\section{Setting}

The germbij note recovers the Taylor jet of a loss at a nondegenerate
minimum from the asymptotic expansions of posterior expectations; the
laplace seabed carries this as the degree-by-degree induction
\code{finite\_jet\_recovery\_of\_monomial\_rates}, whose data at step $k$
are the $d^k$ monomial words $x \mapsto \prod_j x_{m_j}$. The user's
question for this tide was which sub-families of observables suffice,
whether there is a clean condition, and what is lost when one takes
fewer. The seabed's answer was implicit: the step-$k$ engine
\code{tendsto\_pairwise\_normalized\_moment\_difference} states that for
\emph{any} continuous polynomially growing test $P$ the rescaled
difference $(M_1(P,q) - M_2(P,q))/q^{k-2}$ tends to
$-\Cov_\gamma[P, Q_k]$, where $Q_k$ is the degree-$k$ Taylor difference
of the two losses. Everything a family can or cannot see at the leading
rate is in that pairing. The tide made the statement explicit.

The deliberation (tide log
\code{tide-log/2026-09-19-14-40-tide-germbij-sufficient.md}, consult in
\code{gpt\_germbij\_sufficient\_v1.md}) proposed the sufficiency theorem
(A), its converse (B), a finite-family obstruction (C), and an instance
construction (B$'$: perturb a certified loss by a kernel polynomial).
Both parties voted A+B+C, with B$'$ deferred as a one-line remark for
the paper.

\section{The thing formalised}

\begin{quote}\itshape
Let $H_k$ be the degree-$k$ homogeneous polynomials on $\R^d$ (the span
of the monomial words). Let $\phi_i$, $i \in \iota$, be continuous with
polynomial growth, and let $\mathcal T_S(Q) = (\Cov_\gamma[\phi_i,
Q])_i$, a linear map $H_k \to \R^\iota$.
\begin{enumerate}
\item[(B)] For certified losses $L_1, L_2$ with equal jets below $k >
2$: all differences $M_1(\phi_i, q) - M_2(\phi_i, q)$ are $o(q^{k-2})$
iff $\mathcal T_S(Q_k) = 0$.
\item[(A)] If $\ker \mathcal T_S = 0$, such data identify the degree-$k$
derivative tensor. The monomial family has $\ker \mathcal T_S = 0$.
\item[(C)] If $d \ge 2$ and $|\iota| = n \le k$, then $\ker
\mathcal T_S \neq 0$; consequently there is a nonzero $Q \in H_k$ such
that any two certified losses with equal lower jets and $Q_k = Q$ have
all their $S$-data $o(q^{k-2})$.
\end{enumerate}
\end{quote}

\noindent The Lean signature of (A):
\begin{verbatim}
theorem iteratedFDeriv_recovery_of_family_rates (hk : 2 < k)
    (A1 : HigherLaplaceDomain k L1 H) (A2 : HigherLaplaceDomain k L2 H)
    (hlower : forall j < k,
      iteratedFDeriv R j L1 0 = iteratedFDeriv R j L2 0)
    (hsymm1 : (iteratedFDeriv R k L1 0).IsSymm)
    (hsymm2 : (iteratedFDeriv R k L2 0).IsSymm)
    {iota : Type*} (phi : iota -> EuclidD d -> R)
    (hphic : forall i, Continuous (phi i))
    (hphig : forall i, HasPolynomialGrowth (phi i))
    (hinj : forall Q in homogPolySpan d k,
      (forall i, gaussianCovariance H (phi i) Q = 0) -> Q = 0)
    (hdata : forall i,
      (fun q => A1.rescaledMoment (phi i) q - A2.rescaledMoment (phi i) q)
        =o[nhdsWithin 0 (Ioi 0)] fun q => q ^ (k - 2)) :
    iteratedFDeriv R k L1 0 = iteratedFDeriv R k L2 0
\end{verbatim}
with
\begin{center}\code{homogPolySpan d k := Submodule.span R (Set.range monomialTest)}.\end{center}
The other declarations:
\begin{itemize}
\item the operator itself, \code{pairingMap}, a linear map on the span;
\item (B): \code{family\_rates\_iff\_pairing\_eq\_zero};
\item (C): \code{exists\_kernel\_of\_finite\_family} and
\code{finite\_family\_leading\_rate\_blind};
\item the monomial instance: \code{monomialTest\_family\_injective}.
\end{itemize}

The hypotheses of (A) and (B) are those of the seabed's step-$k$
theorem; nothing was added. The injectivity hypothesis quantifies over
the finite-dimensional span, not over all continuous homogeneous
functions, on Astra's advice: the latter space is infinite-dimensional
for $d \ge 2$ and no finite family could ever satisfy it, so the theorem
would have been true and useless.

\section{Proof strategy}

(B) is the pairing limit plus uniqueness of limits in one direction, and
the pairing limit with target $0$ plus \code{isLittleO\_iff\_tendsto'} in
the other. (A) is (B) followed by injectivity, then the diagonal
argument copied from \code{DegreeRecovery}: $Q_k = 0$ gives equal
diagonals of the two symmetric tensors, and
\code{iteratedFDeriv\_eq\_of\_diag\_eq} polarises. The one new fact is
that $Q_k$ lies in the span, which is the seabed's own
\code{diag\_eq\_sum\_monomialTest} read as a membership statement.

For (C) the span is finite-dimensional
(\code{FiniteDimensional.span\_of\_finite}) and the $k + 1$ two-coordinate
words $x_0^{k-j} x_1^{j}$ are linearly independent: evaluate any
relation along the line $(s, 1, 0, \dots)$ to get a real polynomial in
$s$ vanishing identically, kill it with \code{Polynomial.funext}, and
read off each coefficient. So $\dim H_k \ge k + 1 > n = \dim \R^n$ and
\code{LinearMap.ker\_ne\_bot\_of\_finrank\_lt} produces the kernel
element. The two-coordinate count is sharp for $d = 2$; the full count
$\binom{k+d-1}{d-1}$ fails earlier in higher dimension and was not
needed. Linearity of the covariance in the second slot, which the
seabed did not have, is three short lemmas from
\code{integrable\_mul\_quadKernel\_of\_polynomialGrowth}.

For the monomial instance, pairings against all words control the
self-pairing of any span element (induction on the first slot with
\code{gaussianCovariance\_comm}), and the seabed's rigidity theorem
finishes.

\section{Roadblocks and resolutions}

Two rounds, all mechanical.

\emph{Section variables versus explicit arguments.} Declaring
\code{variable (hd : 2 <= d)} and also writing \code{(hd : 2 <= d)} on
the definitions double-bound the name, so every later application
\code{sliceWord hd j} fed \code{j} into the proposition slot. Dropping
the section variable fixed twelve errors at once.

\emph{\code{rw} with lemmas whose function arguments are lambdas.} The
rewrite
\begin{center}\code{rw [<- gaussianExpectation\_add hH (hphic.mul hPc) ...]}\end{center}
instantiated the lemma from the continuity proofs, whose statement
prints the product as \code{phi * P}, and then failed to find that
pattern. Elaborating the lemma first with explicit \code{(f := fun x =>
...)} and closing with \code{Eq.trans} and \code{congr 1; funext} avoids
the pattern match entirely. Same family as the \code{Pi.add} gotcha in
\code{CLAUDE.md}.

\emph{\code{split\_ifs} did not split an \code{if} inside a coordinate
projection.} A \code{by\_cases} followed by two \code{rw [if\_pos hi]}
did. And on this pin \code{WithLp} is a structure, so
\code{EuclideanSpace.single\_apply} is deprecated for
\code{PiLp.single\_apply}, and \code{simp [Fin.ext\_iff]} alone decides
the resulting \code{if}s.

\emph{The consult.} One round, before writing, and it changed the
statements rather than the proofs. The load-bearing paragraph on my
working note's claim that the pairing depends only on the jet of the
observable:
\gptbox{Already in one dimension, set
\[\phi(x) = x^k - c\,x^{k+2}, \qquad c = \Cov(X^k, X^k)/\Cov(X^{k+2}, X^k).\]
Then $\phi$ has the same
$k$-jet as $x^k$, but $\Cov(\phi(X), X^k) = 0$. A genuinely $H$-uniform
sufficient condition is $\mathcal H_k \subseteq \operatorname{span}\{\phi_i\}
+ \R\mathbf 1$ as functions, not merely as jets.}
And on the scope of (C):
\gptbox{C+B+B$'$ do not prove that no finite family recovers the full
jet from its entire expansions. A perturbation invisible at order
$q^{k-2}$ may be visible later. For example, with $H = I$, take locally
$L_\varepsilon(x) = \tfrac12|x|^2 + \varepsilon x_1^3$, $\phi(x) =
x_1^2$. The degree-three covariance vanishes by parity. Nevertheless
$M_\varepsilon(q) - M_0(q) = 45\varepsilon^2 q^2 + o(q^2)$.}
The coefficient checks by hand ($\mathbb E X^8 - \mathbb E X^2\,
\mathbb E X^6 = 105 - 15 = 90$, halved). Both corrections are now in
the file's module docstring and in the working notes for the paper.

\section{What was Mathlib, what was new}

Mathlib, all lookup:
\begin{itemize}
\item spans: \code{Submodule.span\_induction} (the current
motive-with-membership form) and \code{FiniteDimensional.span\_of\_finite};
\item dimension: \code{LinearIndependent.of\_comp} and
\code{fintype\_card\_le\_finrank};
\item kernel: \code{Module.finrank\_fin\_fun} and
\code{LinearMap.ker\_ne\_bot\_of\_finrank\_lt};
\item polynomials: \code{Polynomial.funext}, \code{finsetSum\_coeff},
and \code{coeff\_C\_mul\_X\_pow};
\item counting: \code{Fin.card\_filter\_val\_lt};
\item asymptotics: \code{isLittleO\_iff\_tendsto'}.
\end{itemize}

Seabed: the pairing limit, the Gaussian rigidity theorem, the diagonal
expansion into monomial words, the certificates of the words, the
polarisation lemma, the integrability of polynomially growing
observables against the quadratic kernel.

New: closure of \code{HasPolynomialGrowth} and
\code{IsHomogeneousOfDegree} under the module operations, the span
\code{homogPolySpan} with its certificates, membership of the Taylor
difference, linearity and symmetry of \code{gaussianCovariance}, the
observation operator, and the six theorems. No new analysis.

\section{Lessons}

\begin{itemize}
\item For \code{CLAUDE.md}: never combine \code{variable (h : P)} with
an explicit \code{(h : P)} on a declaration in the same section; and
elaborate lambda-argument lemmas with named \code{(f := ...)} and close
by \code{Eq.trans} rather than \code{rw}.
\item For the tide skill: when a statement is about ``which data
suffice'', consult before writing on the \emph{quantifier domain} of the
condition. The move from all homogeneous functions to the polynomial
span cost nothing in Lean and is the difference between a theorem
finite families can meet and one they cannot.
\item For the paper: a rate-level identifiability theorem needs its
scope stated as such. The working notes had drifted from ``degree-$D$
observables give the $D$-jet'' to ``and nothing more'', which the
seabed does not prove and which the cubic example shows is false at the
level of full expansions.
\end{itemize}

\section{Follow-ups}

\begin{itemize}
\item \emph{Instance (B$'$).} From a kernel direction $Q$, the pair
$L_0 = \tfrac12 H(x,x)$, $L_1 = L_0 + Q$ realises the blindness with
certified packages: both are polynomials, so the order-$k$ Taylor
remainder of $L_1$ is zero and the quadratic lower bound survives a
shrunken ball. Needs the Taylor coefficients of a polynomial in the
\code{taylorHomogeneousTerm} normalisation.
\item \emph{Recovery modulo the invisible subspace.} The data determine
the $\Cov_\gamma$-orthogonal projection of $Q_k$ onto $(\ker
\mathcal T_S)^\perp$; the seabed's rigidity theorem makes $\Cov_\gamma$
an inner product on $H_k$. The cheap Lean form is
$\mathcal T_S Q = \mathcal T_S R \iff Q - R \in \ker \mathcal T_S$; the
substantive one is a quantitative singular-value bound for stable
recovery from noisy leading coefficients.
\item \emph{Full-expansion identifiability from a finite family.} Open.
The cubic example suggests a finite family may recover the jet from all
orders through nonlinear terms in the earlier perturbations; the seabed
has the forward all-orders expansions but no machinery for the
second-order coefficient in the perturbation.
\item \emph{The $H$-uniform spanning criterion.} $H_k \subseteq
\operatorname{span}\{\phi_i\} + \R\mathbf 1$ as functions implies
injectivity for every positive-definite $H$; a short corollary of the
monomial instance and covariance linearity.
\end{itemize}

\section*{Acknowledgements}

GPT-6 Astra corrected the quantifier domain of the sufficiency theorem
and the scope of the obstruction, and supplied both counterexamples.
The seabed is the germbij nondegenerate arc, in particular the
monomial-tests tide whose diagonal expansion this file reuses, and the
normalized-singular tide of the same day,
\tideref[2026-09-19-06-05]{2026-09-19-06-05-tide-germbij-normalized},
whose retrospective sets the paper context.

\end{document}
